\documentclass[11pt]{article}

% Load the 'helvet' package to set Helvetica as the default sans-serif font
\usepackage{helvet}

\usepackage{enumitem}
\setlist{nosep} % reduce space between items globally

% Set the default font family to sans-serif
\renewcommand{\familydefault}{\sfdefault}

% Load the 'geometry' package to adjust page margins
\usepackage[a4paper, margin=1in]{geometry}

\begin{document}

\title{GFET-modulated enzymatic DNA synthesis}

\author{Erik Garrison}
%\date{April 2023}
\maketitle

%\today

\section*{Introduction}

This is a description of a set of protocols that can be run on a graphene field effect transistor or GFET system.
We will be investigating if we can influence the behavior of a transferase enzyme that attaches new nucleotides to the 3' end of a single-stranded DNA molecule.
These terminal deoxynucleotidyl transferase (TdT) enzymes are used to achieve very low synthesis costs, but they are promiscuous and can't easily be controlled or made specific in their activity, so they are usually engaged in the context of a microfluidic platform.

We believe that by changing the effective isolectric point of single-stranded DNA bound to the GFET, we should be able to block interaction between the TdT and the DNA sequence.
In theory, this would allow for sequence-specific patterning of single DNA molecules on an array of GFETs exposed to a flow of TdT+dNTPs.
To test this, we want to build a simple system that focuses on synthesis of a single pattern on a collection of GFETs.

\subsection*{System setup}

Our target system for testing is the S20 or S21 from Graphenea.
These are graphene field effect transistor chips that have a number of distinct graphene channels.
Both chips are modulated by a single liquid gate electrode.
These are either linked into the same source and drain connections (S21) or have separate sensing of pairs of GFETs (S20).
We plan to modify the Graphenea cartridge system which is used to hold these chips during biological sensing applications.
The Graphenea cartridge allows for liquid measurements using these passivated GFETs and includes a breakout board with connectors for attachment to source/measure devices.

Hypothetically, we would replace the liquid manifold that they have, which allows only input of a sample, with one that is capable of inducing a flow across the graphene field effect transistor channels.
The objective of this is to enable tests of serial application of different reactants to the GFET channels under different electrostatic configurations.
The basic idea is that enzymatic DNA synthesis can be modulated, or that is turned on and off using different electrical states at the GFET channel-liquid interface.
In principle, this could be modulated by changing the liquid gate voltage.
This would affect all of the GFETs on these chips simultaneously.
Potentially, we could also envision modifying channel voltages in order to invoke the same kind of result.

\subsection*{Considerations}

The liquid gate voltage range is limited by the voltage of water electrolysis, which occurs at around 1.5 volts.
In later practice, we may use a back-gated GFET system, but there are no commercial platforms which are suitable.
The goal with this initial work is to test the feasibility of guiding enzymatic DNA synthesis using electrical inputs.
This means that we will be simulating the kinds of effects that we would see on single GFETs in a parallel setting under a macroscopic flow or flood context.
We would like to measure if enzymatic synthesis can be blocked by setting a particular global electrostatic state, and also for what states enzymatic synthesis can proceed.
We might also attempt to enable the incorporation of only certain DNA bases out of a flow that may include a series of different DNA bases as reagents that flow across the surface of the chip.

\section*{Protocol 1: Can enzymatic DNA synthesis be modulated?}

Our first protocol tests if enzymatic DNA synthesis can be modulated at all.
We would first functionalize the graphene channel surface with single-stranded DNA template molecules.
Subsequently, we would fix a specific voltage on either the liquid gate or on the channels themselves.
We would then flow across reagents, perhaps a series of TdT+dNTPs of alternating DNA bases.
Finally, we would apply an enzymatic cleavage step that would target a sequence found in the single-stranded DNA templates to cleave the created molecules from the graphene channels.
We would collect this particular wash for analysis with capillary electrophoresis or nanopore sequencing.
The output would be the length distribution of the synthesized DNA molecules.

If we can completely block synthesis, then we would expect to be able to observe this here, as all molecules would be just the fragment of the DNA template molecules downstream of the restriction site.
We should also be able to obtain a positive control that enzymatic synthesis can progress by running the system in a neutral electrical configuration.

\subsection*{Protocol 1 Steps}
\begin{enumerate}
    \item Start by functionalizing the graphene channel surface with single-stranded DNA template molecules.
    \begin{enumerate} %[i.]
        \item Functionalize the graphene channel surface with pyrenebutyric acid N-hydroxysuccinimide ester (PBASE) via $\pi$-$\pi$ stacking.
        \item Introduce single-stranded DNA molecules that have been modified to contain a single amine group at the 5’ end. The amine group forms an amide bond with the carboxyl group of a surface-functionalized PBASE molecule, resulting in stably end-tethered DNA molecules.
        \item Rinse the chip with a buffer solution to wash away non-specific, weakly adsorbed DNA molecules.
    \end{enumerate}
    \item Set a specific voltage on either the liquid gate or on the channels themselves using the source measure unit (SMU). Start with a lower voltage and gradually increase, staying within the safe range of up to 1.2V.
    \item Flow an alternating sequence of A, T, G, C across the chip surface. Use approximately 50µL of each base and ensure that the channels leading up to the chip surface are properly flushed.
    \item Incubate the system for a duration that optimizes the enzymatic synthesis. The time period could vary between 1 to 10 minutes.
    \item Inactivate and remove the terminal deoxynucleotidyl transferase (TdT) before release, so that it cannot continue extending the released products in the collected wash.
    \begin{enumerate} %[i.]
        \item Wash the channels with buffer to remove residual TdT and dNTPs.
        \item Heat-inactivate at 70--75 °C for 10 minutes (or quench with EDTA), then exchange into HindIII reaction buffer (Mg$^{2+}$, Co$^{2+}$-free).
    \end{enumerate}
    \item Form the duplex required for cleavage and cleave the products from the surface using HindIII.
    \begin{enumerate} %[i.]
        \item Flow the 5'-fluoroscein-labelled reverse complement \texttt{fluor-tgtaatc} (50\,nt) over the chip and anneal with a temperature ramp: heat to about 90 °C, then cool slowly to 37 °C. The full-length duplex melts at high temperature, so a 37 °C soak alone will neither denature nor reanneal it.
        \item Flow HindIII and incubate at 37 °C to cleave the synthesized products from the surface of the GFET.
    \end{enumerate}
    \item Collect the wash from this final step (approximately 50-100µL) and analyze it using capillary electrophoresis or nanopore sequencing to observe the length distribution of the synthesized DNA molecules. The cut is internal (between nt 13 and 14 of the template), so the no-synthesis control is a 37 nt product and synthesis shifts the distribution to 37 + $n$ nt. Because the fluorophore sits on the complement strand, denaturing capillary electrophoresis reports how much was released rather than product length; use nanopore sequencing for length distributions.
\end{enumerate}

\section*{Protocol 2: Specifying the order of incorporated DNA bases}

If we can obtain positive results on the previous test, then we might think about implementing specific patterns of DNA bases that are a subset the total ordered flow of DNA bases.
For instance, our flow might be A, T, G, C, and so on with a wash between each one.
But if we can configure the GFETs electrically to prevent synthesis during all but one of the bases in each set, we should be able to write specific sequences.
For example, even if the flow has all bases in it, we could in theory create DNA sequences that just had A homopolymers or T homopolymers, or which were GCGCGC repeats, and so on.
In effect, this protocol is very similar to the previous one, but instead of just having perhaps a single DNA base or single setting of the electrical state on the GFETs, we would be modulating that to match the flow in particular that's being made.

The overall process is the same as in Protocol 1.
We would proceed until such a point that we decided to stop the reaction and at that point we would use an enzyme to cleave the synthesized products from the surface of the GFET.
We would collect the wash from the final step and then analyze that to see how well we did.
In this case we might apply both capillary electrophoresis and DNA sequencing to be sure that we've synthesized the correct molecules.

\subsection*{Protocol 2 Steps}
\begin{enumerate}
    \item Follow steps 1-2 from Protocol 1.
    \item Flow a cyclic sequence of A, T, G, C across the chip surface. Use approximately 50µL of each base.
    \item Modulate the electrical state on the GFETs to match the flow using a programmable source measure unit (SMU). The state should prevent synthesis during all but one of the bases in each set.
    \item Incubate the system for a duration that optimizes the enzymatic synthesis. The time period could vary between 1 to 10 minutes.
    \item After incubation, inactivate and remove TdT, anneal the complement and cleave with HindIII, following step 5 of Protocol 1.
    \item Collect the wash from this final step (approximately 50-100µL). Analyze the wash using capillary electrophoresis and DNA sequencing to confirm the intended sequence pattern.
\end{enumerate}

\section*{Release chemistry: duplex formation and internal HindIII cleavage}

The template oligonucleotide is single-stranded when it is tethered to the graphene and throughout TdT extension, because TdT requires a single-stranded 3' hydroxyl to extend.
The release step therefore has to create a short duplex around the restriction site on demand.
We use HindIII, which cuts the internal site AAGCTT at positions 13--18 of the 50\,nt template and is inactive on single-stranded DNA, and therefore requires that a complementary strand be annealed before digestion.
We use the 5'-fluoroscein-labelled 50\,nt oligo we already hold (\texttt{fluor-tgtaatc}, primer 515677F09, order 12091076), which is the exact reverse complement of the template.
We deliberately do not use the Type IIS sites present in the template (Esp3I/BsmBI, Eco31I/BsaI) for the release step: their reported activity on single-stranded DNA is none (BsaI) or at best partial (BsmBI), and FokI requires a duplex recognition site together with dimerization to cleave.
They remain useful for nested cuts made \emph{in solution} on the released product once it is duplexed.

Because the 5' end of the template is nt 1 and the HindIII cut falls between nt 13 and 14, the released molecule is the 37\,nt downstream fragment plus whatever TdT has appended, with a 4\,nt 5'-AGCT overhang from the staggered cut.
The no-synthesis control is therefore a 37\,nt species, and synthesis moves the product to 37 + $n$\,nt.

\texttt{fluor-tgtaatc}: 5'-TGAGACCAGAGACGAGTTAAGACCTACATCCTAAGCTTGCCTGTGTAATC-3', 50\,nt, 5'-fluoroscein (primer 515677F09, order 12091076, 19.7\,nmol).
The anneal is maintained as a dated bench protocol outside this document: \texttt{protocols/2026-11-11-fluor-tgtaatc-hybridization.md}, which also covers the fluorescent readout of a hybridized chip and the one-way consequence of this step.
Its stock is annotated 100\,$\mu$M by hand with no reconstitution volume recorded, so remeasure it before use rather than trusting the label.
As a full-length complement it forms a blunt 50\,bp duplex with the template, which duplexes the 3' end and blocks TdT, so it must be annealed at the release step only and never before or during extension.
By contrast, a short complement spanning only the restriction site could be pre-annealed, but would then need a blocked 3' end (3'-dideoxy, 3'-phosphate or inverted dT), because a free 3' end is itself a TdT substrate.

HindIII cuts this duplex at the same position, so the complement is split too: the 33\,nt portion paired to template nt 18--50 stays with the released product and carries the fluorophore, while the 17\,nt portion paired to nt 1--17 remains on the surface stub.
The released species is therefore the 37\,nt template fragment duplexed over 33\,bp, with a 4\,nt 5'-AGCT overhang at its 5' end and a single-stranded 3' tail of $n$ nucleotides.

Note what this means for readout.
The fluorescent strand is the complement, and its length is fixed at 33\,nt however much TdT added.
Under denaturing capillary electrophoresis the label therefore reports \emph{how much} was released, not \emph{how long} the products are.
Length distributions require nanopore sequencing of the released material, or non-denaturing CE, where mobility does depend on total length.
The fluorophore is best used for release efficiency and quantitation, and sequencing for length.

Two further requirements follow from this design.
First, TdT must be inactivated and washed out before release, otherwise it continues to extend the released fragments in the collected wash and smears the length distribution; note that this also makes a failed release indistinguishable from blocked synthesis unless a release control is included.
Second, the TdT reaction is normally run in Co$^{2+}$ buffer, which must be exchanged for a Mg$^{2+}$ HindIII buffer, since Co$^{2+}$ inhibits HindIII and can relax its specificity.

\section*{Reagents Required for Protocols}

We require at least 10 separate reagents for these protocols (note the four types of dNTPs).
For flexibility, the system should contain 16 reservoirs.

\begin{enumerate}
    \item Single-stranded DNA template molecules with a 5' amine group and an internal HindIII recognition site (AAGCTT at nt 13--18 of the 50 nt template).
    \item Pyrenebutyric acid N-hydroxysuccinimide ester (PBASE).
    \item Buffer solution for rinsing the chip and for dilution of reagents as required.
    \item Deoxynucleotide triphosphates (dNTPs): dATP, dTTP, dGTP, and dCTP.
    \item Terminal deoxynucleotidyl transferase (TdT).
    \item HindIII restriction enzyme, with 10x HindIII reaction buffer (Mg$^{2+}$, Co$^{2+}$-free).
    \item Single-stranded DNA primer: the 5'-fluoroscein-labelled reverse complement of the template (\texttt{fluor-tgtaatc}: 5'-TGAGACCAGAGACGAGTTAAGACCTACATCCTAAGCTTGCCTGTGTAATC-3'), which forms the duplex that HindIII requires.
\end{enumerate}

\section*{Graphene Functionalization Protocol}

The functionalization procedure is maintained as dated bench protocols outside this document: \texttt{protocols/2026-11-04-pbase-dmf-graphene.md} for PBASE coating, and \texttt{protocols/2026-11-05-amino-gattaca-pbase-graphene.md} for attachment of the amino-modified oligo and ethanolamine blocking.
Those are the procedures we run: about 5\,mM PBASE in dry DMF for 100\,min with the surface kept wet, then 1\,$\mu$M 5'-amine DNA for 60\,min at room temperature, then 100\,mM ethanolamine at pH 9.0 for 30\,min.
An earlier version of this document carried an inline procedure (1\,mM PBASE dried by solvent evaporation for 15--30\,min, 50\,mM ethanolamine in PBS) that was never executed.
It has been removed rather than left in place to be followed by mistake, and the earlier text remains in the repository history.
DMF handling, PPE and waste-disposal requirements are recorded with the coating protocol.

\end{document}
